\documentclass[
    language=english,
    doctype=technical-report,
    institution=none,
    titlestyle=book,
    oneside,
]{../../omnilatex}

\RequirePackage{config/document-settings}

\addbibresource{../../bib/bibliography.bib}

\begin{document}

\maketitle

\tableofcontents

%=============================================================================
\chapter{Introduction}
\label{sec:introduction}
%=============================================================================

This test plan describes the quality assurance strategy for DataPipe
version 3.2. It defines the testing scope, approach, test cases,
environment requirements, and entry and exit criteria. The plan covers
functional, performance, security, and integration testing.

\section{Scope}

This test plan covers:

\begin{itemize}
    \item Core pipeline engine and execution framework.
    \item Kafka and PostgreSQL connectors.
    \item Built-in transforms (filter, map, aggregate, window).
    \item CLI and REST API.
    \item Web dashboard (basic smoke tests).
    \item Exactly-once delivery semantics.
\end{itemize}

\section{Objectives}

\begin{enumerate}
    \item Verify all user stories in the 3.2 release are correctly
        implemented.
    \item Ensure backward compatibility with 3.1 pipeline configurations.
    \item Validate performance targets are met under production-representative
        load.
    \item Confirm no regressions from the 3.1 release.
\end{enumerate}

%=============================================================================
\chapter{Test Strategy}
\label{sec:strategy}
%=============================================================================

\section{Testing Levels}

\begin{table}[htbp]
    \centering
    \caption{Testing levels and responsibilities.}
    \label{tab:test-levels}
    \begin{tabular}{@{}llp{5cm}@{}}
        \toprule
        Level & Owner & Scope \\
        \midrule
        Unit tests & Developers & Individual methods and classes. \\
        Integration tests & Developers & Component interactions
            (connector + pipeline). \\
        System tests & QA team & End-to-end pipeline execution. \\
        Performance tests & SRE team & Throughput and latency under load. \\
        Security tests & Security team & Vulnerability scanning and
            penetration testing. \\
        \bottomrule
    \end{tabular}
\end{table}

\section{Test Types}

\subsection{Functional Testing}

Functional tests verify that each feature behaves according to the
specification. Tests are written in Java using JUnit 5 and executed as
part of the CI pipeline.

\subsection{Performance Testing}

Performance tests use JMH for micro-benchmarks and a custom load
generator for end-to-end throughput measurements. The target is
500\,000 events per second on a three-node cluster.

\subsection{Regression Testing}

The full regression suite runs nightly against the \texttt{main} branch.
Any failure blocks the next release candidate.

%=============================================================================
\chapter{Test Cases}
\label{sec:test-cases}
%=============================================================================

\section{Functional Test Cases}

\begin{table}[htbp]
    \centering
    \caption{Functional test cases for the pipeline engine.}
    \label{tab:func-tests}
    \begin{tabular}{@{}p{2.5cm}p{5cm}p{3cm}p{2cm}@{}}
        \toprule
        Test ID & Description & Expected Result & Priority \\
        \midrule
        TC-001 & Start a pipeline with Kafka source and
            PostgreSQL sink & Records flow from source to sink
            within 10s & P0 \\
        TC-002 & Stop a running pipeline gracefully &
            Pipeline drains in-flight records before stopping & P0 \\
        TC-003 & Apply filter transform & Only matching records
            reach the sink & P0 \\
        TC-004 & Apply windowed aggregation & Aggregate results
            match expected values & P0 \\
        TC-005 & Restart a failed pipeline & Pipeline resumes from
            last checkpoint & P1 \\
        TC-006 & Pipeline with exactly-once delivery & No duplicates
            after consumer restart & P0 \\
        TC-007 & Handle null records in transform & Null records are
            skipped without error & P1 \\
        TC-008 & CLI pipeline list command & All registered pipelines
            are displayed & P1 \\
        TC-009 & REST API pipeline status & JSON response matches
            expected schema & P1 \\
        TC-010 & Pipeline validation with missing config & Clear error
            message identifying the missing field & P1 \\
        \bottomrule
    \end{tabular}
\end{table}

\section{Performance Test Cases}

\begin{table}[htbp]
    \centering
    \caption{Performance test cases.}
    \label{tab:perf-tests}
    \begin{tabular}{@{}p{2.5cm}p{4cm}p{3.5cm}p{2cm}@{}}
        \toprule
        Test ID & Description & Target & Priority \\
        \midrule
        PT-001 & Single pipeline throughput &
            $\geq 500$k events/s on 3 nodes & P0 \\
        PT-002 & End-to-end latency (P99) & $< 200$\,ms & P0 \\
        PT-003 & Pipeline startup time & $< 5$\,seconds & P1 \\
        PT-004 & Consumer lag recovery & Recover from 100k lag
            in $< 60$\,s & P1 \\
        PT-005 & Memory usage under load & $< 8$\,GB per node &
            P1 \\
        PT-006 & Concurrent pipeline execution &
            10 pipelines simultaneously with $< 10\%$ degradation & P2 \\
        \bottomrule
    \end{tabular}
\end{table}

\section{Security Test Cases}

\begin{table}[htbp]
    \centering
    \caption{Security test cases.}
    \label{tab:sec-tests}
    \begin{tabular}{@{}p{2.5cm}p{6cm}p{3cm}p{1.5cm}@{}}
        \toprule
        Test ID & Description & Expected Result & Priority \\
        \midrule
        ST-001 & SQL injection via pipeline config &
            Input is sanitised, no SQL execution & P0 \\
        ST-002 & Unauthenticated API access &
            401 Unauthorized returned & P0 \\
        ST-003 & TLS termination on API endpoints &
            All traffic encrypted in transit & P0 \\
        ST-004 & Secret rotation for connector credentials &
            Pipeline reloads new credentials without restart & P1 \\
        \bottomrule
    \end{tabular}
\end{table}

%=============================================================================
\chapter{Test Environment}
\label{sec:environment}
%=============================================================================

\section{Hardware Requirements}

\begin{table}[htbp]
    \centering
    \caption{Test environment hardware specifications.}
    \label{tab:hardware}
    \begin{tabular}{@{}lll@{}}
        \toprule
        Component & Specification & Quantity \\
        \midrule
        Application nodes & 8 vCPU, 32\,GB RAM, 200\,GB SSD & 3 \\
        Kafka brokers & 4 vCPU, 16\,GB RAM, 500\,GB SSD & 3 \\
        PostgreSQL & 4 vCPU, 16\,GB RAM, 200\,GB SSD & 2 (primary + replica) \\
        Load generator & 8 vCPU, 16\,GB RAM & 1 \\
        \bottomrule
    \end{tabular}
\end{table}

\section{Software Requirements}

\begin{itemize}
    \item Ubuntu 22.04 LTS on all nodes.
    \item Java 21 (Eclipse Temurin).
    \item Docker 24+ and Docker Compose v2.
    \item Kafka 3.7 with KRaft mode.
    \item PostgreSQL 16.
    \item DataPipe 3.2.0-rc1 (release candidate).
\end{itemize}

\section{Test Data}

\begin{itemize}
    \item Synthetic dataset: 10 million records in JSON format.
    \item Production-representative dataset: anonymised snapshot from
        staging environment.
    \item Edge-case dataset: records with null fields, unicode characters,
        and maximum-size payloads.
\end{itemize}

%=============================================================================
\chapter{Entry and Exit Criteria}
\label{sec:criteria}
%=============================================================================

\section{Entry Criteria}

Testing may begin when:

\begin{enumerate}
    \item All P0 features are code-complete.
    \item Unit test coverage meets the minimum threshold (80\%).
    \item The release candidate build passes the CI pipeline.
    \item Test environment is provisioned and verified.
    \item Test data is loaded and validated.
\end{enumerate}

\section{Exit Criteria}

Testing is complete when:

\begin{enumerate}
    \item All P0 test cases pass.
    \item 95\% of P1 test cases pass.
    \item No open P0 or P1 defects.
    \item Performance targets are met.
    \item Security scan reports no critical vulnerabilities.
    \item Regression suite is clean.
    \item Release notes are drafted and reviewed.
\end{enumerate}

\section{Defect Severity Definitions}

\begin{table}[htbp]
    \centering
    \caption{Defect severity levels and response targets.}
    \label{tab:defects}
    \begin{tabular}{@{}lp{5cm}p{3cm}p{2cm}@{}}
        \toprule
        Severity & Description & Fix Target & Blocks Release? \\
        \midrule
        S1 & Data loss or corruption & 24 hours & Yes \\
        S2 & Feature unusable, no workaround & 48 hours & Yes \\
        S3 & Feature degraded, workaround available & Next sprint & No \\
        S4 & Cosmetic or minor issue & Backlog & No \\
        \bottomrule
    \end{tabular}
\end{table}

\printbibliography[heading=bibintoc, title={References}]

\end{document}
